% ZHOU-SOURCE-PAGE: 49 PRINTED: 33
\section{Calculus and Linear Algebra}

Calculus and linear algebra lay the foundation for many advanced math topics used in
quantitative finance. So be prepared to answer some calculus or linear algebra
problems---many of them may be incorporated into more complex problems---in
quantitative interviews. Since most of the tested calculus and linear algebra knowledge
is easy to grasp, the marginal benefit far outweighs the time you spend brushing up your
knowledge on key subjects. If your memory of calculus or linear algebra is a little rusty,
spend some time reviewing your college textbooks!

Needless to say, it is extremely difficult to condense any calculus/linear algebra books
into one chapter. Neither is it my intention to do so. This chapter focuses only on some
of the core concepts of calculus/linear algebra that are frequently occurring in
quantitative interviews. And unless necessary, it does so without covering the proof,
details or even caveats of these concepts. If you are not familiar with any of the concepts,
please refer to your favorite calculus/linear algebra books for details.

\subsection{Limits and Derivatives}

\subsubsection*{Basics of derivatives}

Let's begin with some basic definitions and equations used in limits and derivatives.
Although the notations may be different, you can find these materials in any calculus
textbook.

\begin{concept}{Derivative}
Let $y=f(x)$, then
\[
f'(x)=\frac{\dd y}{\dd x}=\lim_{\Delta x\to0}\frac{\Delta y}{\Delta x}
=\lim_{\Delta x\to0}\frac{f(x+\Delta x)-f(x)}{\Delta x}.
\]
\end{concept}

\begin{concept}{The product rule}
If $u=u(x)$ and $v=v(x)$ and their respective derivatives exist,
\[
\frac{\dd(uv)}{\dd x}=u\frac{\dd v}{\dd x}+v\frac{\dd u}{\dd x},
\qquad (uv)'=u'v+uv'.
\]
\end{concept}

\begin{concept}{The quotient rule}
\[
\frac{\dd}{\dd x}\left(\frac{u}{v}\right)
=\frac{v\,\dfrac{\dd u}{\dd x}-u\,\dfrac{\dd v}{\dd x}}{v^2},
\qquad
\left(\frac{u}{v}\right)'=\frac{u'v-uv'}{v^2}.
\]
\end{concept}

\begin{concept}{The chain rule}
If $y=f(u(x))$ and $u=u(x)$, then
\[
\frac{\dd y}{\dd x}=\frac{\dd y}{\dd u}\frac{\dd u}{\dd x}.
\]
\end{concept}

\begin{concept}{The generalized power rule}
\[
\frac{\dd y^n}{\dd x}=ny^{n-1}\frac{\dd y}{\dd x}\qquad\text{for }\forall n\ne0.
\]
\end{concept}

\begin{concept}{Some useful equations}

% ZHOU-SOURCE-PAGE: 50 PRINTED: 34
\[
a^x=e^{x\ln a},
\qquad \ln(ab)=\ln a+\ln b,
\qquad e^x=\lim_{n\to\infty}\left(1+\frac{x}{n}\right)^n.
\]

\[
\lim_{x\to0}\frac{\sin x}{x}=1,
\qquad
\lim_{x\to0}(1+x)^k=1+kx\quad\text{for any }k.
\]

\[
\lim_{x\to\infty}\frac{\ln x}{x^r}=0\quad\text{for any }r>0,
\qquad
\lim_{x\to\infty}x^r e^{-x}=0\quad\text{for any }r.
\]

\[
\frac{\dd}{\dd x}e^u=e^u\frac{\dd u}{\dd x},
\qquad
\frac{\dd}{\dd x}a^u=(a^u\ln a)\frac{\dd u}{\dd x},
\qquad
\frac{\dd}{\dd x}\ln u=\frac{1}{u}\frac{\dd u}{\dd x}=\frac{u'}{u}.
\]

\[
\frac{\dd}{\dd x}\sin x=\cos x,
\qquad
\frac{\dd}{\dd x}\cos x=-\sin x,
\qquad
\frac{\dd}{\dd x}\tan x=\sec^2 x.
\]
\end{concept}

\begin{problem}{What is the derivative of $y=(\ln x)^{\ln x}$?}
\hint{To calculate the derivative of functions with the format $y=f(x)^x$, it is common to take natural logs on both sides and then take the derivative, since $\dd(\ln y)/\dd x=1/y\cdot\dd y/\dd x$.}
\end{problem}

\begin{solution}

This is a good problem to test your knowledge of basic derivative formulas---
specifically, the chain rule and the product rule.

Let $u=\ln y=\ln\bigl((\ln x)^{\ln x}\bigr)=\ln x\cdot\ln(\ln x)$. Applying the chain rule and the product rule,
we have
\[
\begin{aligned}
\frac{\dd u}{\dd x}&=\frac{\dd(\ln y)}{\dd x}=\frac{1}{y}\frac{\dd y}{\dd x}\\
&=\frac{\dd(\ln x)}{\dd x}\cdot\ln(\ln x)+\ln x\cdot\frac{\dd(\ln(\ln x))}{\dd x}\\
&=\frac{\ln(\ln x)}{x}+\frac{\ln x}{x\ln x},
\end{aligned}
\]

To derive $\dfrac{\dd(\ln(\ln x))}{\dd x}$, we again use the chain rule by setting $v=\ln x$:
\[
\begin{aligned}
\frac{\dd(\ln(\ln x))}{\dd x}&=\frac{\dd(\ln v)}{\dd v}\frac{\dd v}{\dd x}\\
&=\frac{1}{v}\cdot\frac{1}{x}=\frac{1}{x\ln x}.
\end{aligned}
\]

\[
\begin{aligned}
\therefore\quad
\frac{1}{y}\frac{\dd y}{\dd x}
&=\frac{\ln(\ln x)}{x}+\frac{\ln x}{x\ln x}\\
&\quad\Rightarrow\quad
\frac{\dd y}{\dd x}=\frac{y}{x}\bigl(\ln(\ln x)+1\bigr)\\
&=\frac{(\ln x)^{\ln x}}{x}\bigl(\ln(\ln x)+1\bigr).
\end{aligned}
\]

\end{solution}

\subsubsection*{Maximum and minimum}

Derivative $f'(x)$ is essentially the slope of the tangent line to the curve $y=f(x)$ and
the instantaneous rate of change (velocity) of $y$ with respect to $x$. At point $x=c$, if

% ZHOU-SOURCE-PAGE: 51 PRINTED: 35
$f'(c)>0$, $f(x)$ is an increasing function at $c$; if $f'(c)<0$, $f(x)$ is a decreasing
function at $c$.

\begin{concept}{Local maximum or minimum}
Suppose that $f(x)$ is differentiable at $c$ and is defined
on an open interval containing $c$. If $f(c)$ is either a local maximum value or a local
minimum value of $f(x)$, then $f'(c)=0$.
\end{concept}

\begin{concept}{Second Derivative test}
Suppose the secondary derivative of $f(x)$, $f''(x)$, is
continuous near $c$. If $f'(c)=0$ and $f''(c)>0$, then $f(x)$ has a local minimum at $c$; if
$f'(c)=0$ and $f''(c)<0$, then $f(x)$ has a local maximum at $c$.
\end{concept}

\begin{problem}{Without calculating the numerical results, can you tell me which number is larger, $e^\pi$ or $\pi^e$?}
\hint{Again, consider taking natural logs on both sides; $\ln a>\ln b\Longleftrightarrow a>b$, since $\ln x$ is a monotonously increasing function.}
\end{problem}

\begin{solution}

Let's take natural logs of $e^\pi$ and $\pi^e$. On the left side we have $\pi\ln e$, on the
right side we have $e\ln\pi$. If $e^\pi>\pi^e$,
\[
\begin{aligned}
e^\pi>\pi^e
&\quad\Longleftrightarrow\quad
\pi\cdot\ln e>e\cdot\ln\pi\\
&\quad\Longleftrightarrow\quad
\frac{\ln e}{e}>\frac{\ln\pi}{\pi}.
\end{aligned}
\]

Is it true? That depends on whether $f(x)=\dfrac{\ln x}{x}$ is an increasing or decreasing function
from $e$ to $\pi$. Taking the derivative of $f(x)$, we have
\[
f'(x)=\frac{1/x\cdot x-\ln x}{x^2}=\frac{1-\ln x}{x^2},
\]
which is less than 0 when $x>e$ ($\ln x>1$). In fact, $f(x)$ has global maximum when
$x=e$ for all $x>0$. So $\dfrac{\ln e}{e}>\dfrac{\ln\pi}{\pi}$ and $e^\pi>\pi^e$.

Alternative approach: If you are familiar with the Taylor's series, which we will discuss
in Section 3.4, you can apply Taylor's series to $e^x$:
\[
e^x=\sum_{n=0}^{\infty}\frac{x^n}{n!}=1+\frac{x}{1!}+\frac{x^2}{2!}+\frac{x^3}{3!}+\cdots.
\]
So $e^x>1+x$, $\forall x>0$. Let $x=\pi/e-1$, then
\[
\frac{e^{\pi/e}}{e}>\frac{\pi}{e}
\quad\Longleftrightarrow\quad
e^{\pi/e}>\pi
\quad\Longleftrightarrow\quad
e^\pi>\pi^e.
\]

\end{solution}

\subsubsection*{L'Hospital's rule}

Suppose that functions $f(x)$ and $g(x)$ are differentiable at $x\to a$ and that
$\lim_{x\to a}g'(a)\ne0$.

Further suppose that $\lim_{x\to a}f(a)=0$ and $\lim_{x\to a}g(a)=0$ or that
$\lim_{x\to a}f(a)\to+\infty$ and $\lim_{x\to a}g(a)$

% ZHOU-SOURCE-PAGE: 52 PRINTED: 36
$\to+\infty$, then
\[
\lim_{x\to a}\frac{f(x)}{g(x)}=\lim_{x\to a}\frac{f'(x)}{g'(x)}.
\]
L'Hospital's rule converts the limit from an indeterminate form to a determinate form.

\begin{problem}{What is the limit of $e^x/x^2$ as $x\to\infty$, and what is the limit of $x^2\ln x$ as $x\to0^+$?}
\end{problem}

\begin{solution}

$\lim_{x\to\infty}\dfrac{e^x}{x^2}$ is a typical example of L'Hospital's rule since
$\lim_{x\to\infty}e^x=\infty$ and $\lim_{x\to\infty}x^2=\infty$. Applying L'Hospital's rule, we have
\[
\begin{aligned}
\lim_{x\to\infty}\frac{f(x)}{g(x)}
&=\lim_{x\to\infty}\frac{e^x}{x^2}\\
&=\lim_{x\to\infty}\frac{f'(x)}{g'(x)}\\
&=\lim_{x\to\infty}\frac{e^x}{2x}.
\end{aligned}
\]

The result still has the property that $\lim_{x\to\infty}f(x)=\lim_{x\to\infty}e^x=\infty$ and
$\lim_{x\to\infty}g(x)=\lim_{x\to\infty}2x=\infty$, so we can apply the L' Hospital's rule again:
\[
\begin{aligned}
\lim_{x\to\infty}\frac{f(x)}{g(x)}
&=\lim_{x\to\infty}\frac{e^x}{x^2}\\
&=\lim_{x\to\infty}\frac{f'(x)}{g'(x)}\\
&=\lim_{x\to\infty}\frac{e^x}{2x}\\
&=\lim_{x\to\infty}\frac{\dd(e^x)/\dd x}{\dd(2x)/\dd x}\\
&=\lim_{x\to\infty}\frac{e^x}{2}=\infty.
\end{aligned}
\]

At first look, L'Hospital's rule does not appear to be applicable to
$\lim_{x\to0^+}x^2\ln x$ since it's not in the format of
$\lim_{x\to a}\dfrac{f(x)}{g(x)}$. However, we can rewrite the original limit as
$\lim_{x\to0^+}\dfrac{\ln x}{x^{-2}}$ and it becomes obvious that
$\lim_{x\to0^+}x^{-2}=\infty$ and $\lim_{x\to0^+}\ln x=-\infty$. So we can now apply
L'Hospital's rule:
\[
\begin{aligned}
\lim_{x\to0^+}x^2\ln x
&=\lim_{x\to0^+}\frac{\ln x}{x^{-2}}\\
&=\lim_{x\to0^+}\frac{\dd(\ln x)/\dd x}{\dd(x^{-2})/\dd x}\\
&=\lim_{x\to0^+}\frac{1/x}{-2/x^3}\\
&=\lim_{x\to0^+}\frac{x^2}{-2}=0.
\end{aligned}
\]

\end{solution}

\subsection{Integration}

\subsubsection*{Basics of integration}

Again, let's begin with some basic definitions and equations used in integration.

\begin{concept}{Antiderivative}
If we can find a function $F(x)$ with derivative $f(x)$, then we call $F(x)$ an
antiderivative of $f(x)$.

If $f(x)=F'(x)$,
\[
\int_a^b f(x)=\int_a^b F(x)\,\dd x=[F(x)]_a^b=F(b)-F(a).
\]

% ZHOU-SOURCE-PAGE: 53 PRINTED: 37
\[
\frac{\dd F(x)}{\dd x}=f(x),\qquad F(a)=y_a\quad\Rightarrow\quad
F(x)=y_a+\int_{a}^{x}f(t)\,\dd t.
\]
\end{concept}

\begin{concept}{The generalized power rule in reverse}
\[
\int u^k\,\dd u=\frac{u^{k+1}}{k+1}+c\qquad(k\ne1),
\]
where $c$ is any constant.
\end{concept}

\begin{concept}{Integration by substitution}
\[
\int f(g(x))\cdot g'(x)\,\dd x=\int f(u)\,\dd u
\quad\text{with }u=g(x),\quad \dd u=g'(x)\,\dd x.
\]
\end{concept}

\begin{concept}{Substitution in definite integrals}
\[
\int_a^b f(g(x))\cdot g'(x)\,\dd x
=\int_{g(a)}^{g(b)}f(u)\,\dd u.
\]
\end{concept}

\begin{concept}{Integration by parts}
\[
\int u\,\dd v=uv-\int v\,\dd u.
\]
\end{concept}

\begin{problem}{A. What is the integral of $\ln(x)$?}
\end{problem}

\begin{solution}

This is an example of integration by parts. Let $u=\ln x$ and $v=x$, we have
\[
\dd(uv)=v\,\dd u+u\,\dd v=(x\cdot1/x)\,\dd x+\ln x\,\dd x,
\]
\[
\therefore\quad\int\ln x\,\dd x=x\ln x-\int\dd x=x\ln x-x+c,
\]
where $c$ is any constant.

\end{solution}

\begin{problem}{B. What is the integral of $\sec(x)$ from $x=0$ to $x=\pi/6$?}
\end{problem}

\begin{solution}

Clearly this problem is directly related to differentiation/integration of
trigonometric functions. Although there are derivative functions for all basic
trigonometric functions, we only need to remember two of them:
\[
\frac{\dd}{\dd x}\sin x=\cos x,\qquad
\frac{\dd}{\dd x}\cos x=-\sin x.
\]
The rest can be derived using the product rule or the quotient rule. For example,
\[
\begin{aligned}
\frac{\dd\sec x}{\dd x}&=\frac{\dd(1/\cos x)}{\dd x}\\
&=\frac{\sin x}{\cos^2 x}=\sec x\tan x,
\end{aligned}
\]
\[
\begin{aligned}
\frac{\dd\tan x}{\dd x}&=\frac{\dd(\sin x/\cos x)}{\dd x}\\
&=\frac{\cos^2 x+\sin^2 x}{\cos^2 x}=\sec^2 x.
\end{aligned}
\]
\[
\therefore\quad
\frac{\dd(\sec x+\tan x)}{\dd x}=\sec x(\sec x+\tan x).
\]
